\section{Introdução}
\label{sec:introducao}

A tecnologia da informação tornou-se parte do funcionamento cotidiano de empresas pequenas: sustenta vendas, atendimento, faturamento, comunicação e continuidade operacional. Para organizações com equipes reduzidas, entretanto, governar essa dependência cria um problema prático. As decisões de TI costumam se concentrar no proprietário e em um profissional interno ou terceirizado, enquanto orçamento, tempo e especialização permanecem limitados. Nesse contexto, a ausência de governança aumenta a exposição a prioridades conflitantes, interrupções e riscos; a adoção literal de estruturas extensas, por outro lado, pode introduzir um custo de coordenação incompatível com a capacidade da organização.

O problema não decorre da falta de referências consolidadas. A ISO/IEC~38500 orienta o corpo governante a avaliar, direcionar e monitorar o uso da TI \citep{iso38500_2024}; o COBIT~2019 oferece componentes e fatores de desenho para sistemas de governança \citep{isaca2018cobit}; o ITIL~4 organiza a gestão de produtos e serviços em torno da criação de valor \citep{axelos2019itil}; e o NIST CSF~2.0 e os CIS Controls~v8 estruturam a gestão de riscos de segurança \citep{nist2024csf,cis2021controls}. O desafio está em transformar esse corpo de conhecimento em um conjunto mínimo, coerente e executável por uma empresa com poucas pessoas, sem descaracterizar os princípios que lhe dão sustentação.

A literatura oferece suporte para uma abordagem contingencial. \citet{levstek2022adaptive} observam que mecanismos de governança de TI não são universalmente necessários e que modelos abrangentes podem ser onerosos para pequenas e médias empresas. Em uma linha complementar, \citet{silva2018baseline} identificam mecanismos estruturais, processuais e relacionais que podem formar uma base mínima de governança, embora seu estudo exploratório tenha uma amostra reduzida. A revisão sistemática de \citet{sagala2024transformation} reforça que a transformação digital em PMEs deve considerar a situação inicial, as limitações e as idiossincrasias da organização, com implantação incremental e aprendizagem contínua. Essas evidências apontam para a necessidade de seleção e adaptação, e não para a simples redução de um único framework.

Este trabalho responde à seguinte questão de pesquisa: \emph{como estruturar um artefato de governança e gestão de TI que preserve mecanismos essenciais de direção, execução e segurança, mas seja operável por empresas com recursos humanos e financeiros restritos?} Como delimitação empírica futura, o estudo concentra sua primeira avaliação em empresas simuladas com 5 a 20 colaboradores. Esse recorte corresponde ao ambiente organizacional para o qual o custo de coordenação e a concentração de papéis tendem a ser mais críticos; ele não autoriza generalizações para todo o universo estatístico das PMEs.

Propõe-se o GEAR (Gestão, Execução, Agilidade e Risco), um artefato contingencial composto por três domínios essenciais: governança e direção, execução e serviços, segurança e continuidade. Uma camada transversal e opcional de assistência por inteligência artificial (IA) apoia diagnóstico, documentação, métricas e consulta ao conhecimento do framework, sempre com revisão humana. A adoção começa por um percurso local de 30 dias, historicamente chamado Fase Zero, destinado a criar visibilidade sobre demandas, decisões, ativos, riscos e indicadores. O prazo orienta o planejamento e não garante implantação ou avanço de maturidade.

O GEAR é apresentado como uma síntese de desenho, não como substituto normativo para ISO/IEC~38500, COBIT, ITIL, NIST CSF ou CIS Controls. O artefato seleciona mecanismos compatíveis com seu contexto-alvo e mantém rastreabilidade até as referências de origem. Essa distinção é especialmente importante na comparação planejada: COBIT trata de governança e gestão empresarial de informação e tecnologia, enquanto ITIL enfatiza gestão de produtos e serviços. Portanto, a avaliação não buscará declarar um vencedor universal, mas examinar o ajuste entre escopo, esforço de adoção e utilidade produzida em um conjunto controlado de tarefas.

A construção do GEAR também coloca uma questão de método: como usar IA na pesquisa, síntese e produção de documentação sem transformar conteúdo gerado em evidência? Este artigo relata o uso ativo dessa assistência com auditoria de fontes, consistência e cálculos, distinguindo a produção do material da aplicação futura do framework. A contribuição é descritiva e rastreável; não se estima um ganho de produtividade da IA nem se presume que a revisão elimine todos os erros.

As contribuições desta etapa são a especificação de um framework ajustável ao contexto de pequenas equipes, o relato da pesquisa e produção assistidas por IA com auditoria de conteúdo e um protocolo para prova de conceito futura. Os 25 casos comparativos da \cref{sec:protocolo} serão executados no FrameSim após o congelamento de cenários e critérios. Tabelas e figuras de resultados permanecem sem valores. Recursos agregados, como MCPs, skills e integrações, são citados como canais associados ao repositório; sua implementação não integra a contribuição ou a validação deste artigo.

O restante do artigo apresenta a literatura e as bases gerenciais na \cref{sec:fundamentacao}, a pesquisa e a auditoria de conteúdo na \cref{sec:metodo}, e a especificação do GEAR na \cref{sec:artefato}. As \cref{sec:estado,sec:protocolo} distinguem o estado documental e a prova de conceito futura. As \cref{sec:discussao,sec:conclusao} discutem o alcance das contribuições e os limites da evidência.
